% =====================================================================
%  DANH MỤC TOÁN 12 — tên bộ tài liệu, chương, chuyên đề.
%  Dùng trong file .tex:  \loai{1}\chuong{3}   |   \loai{2}\de{01}   |   \loai{3}\chuyende{1}
% =====================================================================
\dmuc{loai1}{CÁC CHUYÊN ĐỀ CƠ BẢN TOÁN 12}
\dmuc{loai2}{ĐỀ KIỂM TRA TOÁN 12}
\dmuc{loai3}{CHUYÊN ĐỀ ÔN THI TỐT NGHIỆP THPT}

% Chương SGK Toán 12 Kết nối tri thức (kiểm lại theo sách đang dạy)
\dmuc{chuong1}{Chương I. Ứng dụng đạo hàm để khảo sát và vẽ đồ thị của hàm số}
\dmuc{chuong2}{Chương II. Vectơ và hệ trục toạ độ trong không gian}
\dmuc{chuong3}{Chương III. Các số đặc trưng đo mức độ phân tán của mẫu số liệu ghép nhóm}
\dmuc{chuong4}{Chương IV. Nguyên hàm và tích phân}
\dmuc{chuong5}{Chương V. Phương pháp toạ độ trong không gian}
\dmuc{chuong6}{Chương VI. Xác suất có điều kiện}
\dmucngan{chuong1}{Chương I. Ứng dụng đạo hàm, khảo sát hàm số}
\dmucngan{chuong3}{Chương III. Số đo độ phân tán của mẫu số liệu ghép nhóm}

% Chuyên đề ôn thi tốt nghiệp THPT (đặt tạm, sửa theo ý)
\dmuc{chuyende1}{Chuyên đề 1. Ứng dụng đạo hàm và khảo sát hàm số}
\dmuc{chuyende2}{Chuyên đề 2. Hàm số mũ, hàm số lôgarit}
\dmuc{chuyende3}{Chuyên đề 3. Nguyên hàm, tích phân và ứng dụng}
\dmuc{chuyende4}{Chuyên đề 4. Hình học không gian}
\dmuc{chuyende5}{Chuyên đề 5. Phương pháp toạ độ trong không gian}
\dmuc{chuyende6}{Chuyên đề 6. Thống kê và xác suất}
